% ==============================================================================
% ANEXOS DEL PROTOCOLO DE INVESTIGACIÓN (RESPALDO)
% ==============================================================================
% Archivo de respaldo generado el 27/09/2026
% Proyecto: "Acceso a los servicios de salud con enfoque intercultural y calidad
%           de atención percibida por los usuarios, Ayacucho 2026"
% Autora: Gresly Lucero Pariona Palomino
% Asesor: Dr. Manglio Aguirre Andrade
% ==============================================================================
% NOTA: Este archivo contiene los 5 anexos completos extraídos del protocolo
% principal (protocolo_cualitativo.tex). Para reincorporarlos, copiar el 
% contenido desde \appendix hasta \end{document} de vuelta al archivo principal.
% ==============================================================================

\appendix
\section*{ANEXOS}
\addcontentsline{toc}{section}{ANEXOS}

% ------------------------------------------------------------------------------
% ANEXO 1: GUÍA DE ENTREVISTA BILINGÜE
% ------------------------------------------------------------------------------
\subsection*{Anexo 1: Instrumento Cualitativo -- Guía de Entrevista en Profundidad Bilingüe}
\addcontentsline{toc}{subsection}{Anexo 1: Guía de Entrevista Bilingüe}
\noindent\textbf{Objetivo:} Comprender las vivencias, significados y percepciones de los usuarios respecto al acceso y la calidad del cuidado recibido en los servicios de salud de Ayacucho.\\[0.2cm]
\noindent\textbf{Protocolo Metodológico Fenomenológico:} En el marco del método fenomenológico-hermenéutico y la reducción fenomenológica (\textit{epojé}), la investigadora iniciará la sesión planteando exclusivamente la \textbf{Pregunta Generadora de Apertura Narrativa}, permitiendo que el participante relate con total libertad, espontaneidad y sin interrupciones el flujo completo de su experiencia vivida desde su comunidad. Las siguientes ocho (08) preguntas norteadoras operarán como \textit{guía de estímulo temático o reactivadores del relato}, utilizándose únicamente como recordatorios en caso de que el participante omita de manera espontánea alguna dimensión clave del acceso o de la calidad durante su narración libre.\\[0.3cm]
\noindent\textbf{I. PREGUNTA GENERADORA / DISPARADORA DE APERTURA NARRATIVA:}\\[0.2cm]
\textbf{En Castellano:}\\
\enquote{Cuénteme con toda confianza y en detalle, desde el momento en que salió de su casa hasta que culminó su atención y regresó a su comunidad, ¿cómo fue toda su vivencia y experiencia personal la última vez que acudió a atenderse en este establecimiento de salud?}\\[0.2cm]
\textbf{En Quechua Chanka:}\\
\enquote{Tukuy sunquykiwan willaykullaway allillamanta, wasikimanta lloqsimusqaykimantapacha samaykusqaykikama hinaspa aylluykiman kutinaykikama, ¿imaynataq karqan llapan puririyniyki hinaspa tarikuyniyki qipallata kay hampi wasiman hampichikuq hamusqaykipi?}\\[0.4cm]
\noindent\textbf{II. PREGUNTAS NORTEADORAS / REACTIVADORES TEMÁTICOS DE RECORDATORIO:}\\
\textit{(Utilizar solo si el participante no abordó espontáneamente la dimensión respectiva)}
\begin{enumerate}
    \item ¿Cómo es el viaje desde su comunidad hasta el centro de salud cuando usted o su familia se enferman? ¿Cuánto tiempo y dinero gasta en llegar? \\
    \textit{(¿Imaynatam llaqtaykimanta hamunki hampina wasiman unqusqa kaspayki? ¿Hayka pachatam purinki, haykatam gastanki?)}
    \item ¿Encuentra siempre las medicinas que le recetan en la farmacia de la posta o tiene que comprarlas afuera con su propio dinero? \\
    \textit{(¿Tarikunchu hampikuna hampina wasipi icha qawapichu rantinki qullqillaykiwan?)}
    \item ¿Qué trámites, madrugadas o dificultades debe realizar para conseguir un cupo de atención médica o de enfermería? \\
    \textit{(¿Ima sasachakuykunatam tarinki tutamanta sayaspa cupota chaskinaykipaq?)}
    \item ¿En qué idioma le habla el personal de salud cuando le atiende? ¿Logra entender claramente las explicaciones sobre su tratamiento en su lengua materna? \\
    \textit{(¿Ima simipim rimapayasunki hampiqkuna? ¿Allintachu intillinki hampina simipi rimapayasuptiyki?)}
    \item ¿Cómo siente que tratan sus costumbres, su vestimenta andina y su forma de ver la enfermedad en este establecimiento de salud? \\
    \textit{(¿Imaynatam qawasunki costumbrikita, pachaykita hampina wasipi?)}
    \item ¿Cuánto tiempo tiene que esperar en la sala antes de que le llamen y cómo es la amabilidad del trato durante esa espera? \\
    \textit{(¿Hayka unaytam tiyanki suyaspa waqyamusunaykikama hinaspa imaynataq chay suyasqaykipi runakuna rimapayasunki?)}
    \item ¿Se siente usted seguro(a) y con confianza al momento de ser examinado(a) por el médico o la enfermera? \\
    \textit{(¿Kallpanchasqachu hinaspa hawkallachu tarikunki doctorkuna o enfermerakuna kurkuykita qawapayasuptiyki?)}
    \item ¿Recomendaría usted a sus familiares y vecinos acudir a este centro de salud cuando tengan un problema de salud? ¿Por qué? \\
    \textit{(¿Nisunkimanchu aylluykikunata kay hampi wasiman hamunankupaq unqusqankupi? ¿Imanasqa?)}
\end{enumerate}

\vspace{0.5cm}

% ------------------------------------------------------------------------------
% ANEXO 2: MATRIZ DE CONSISTENCIA
% ------------------------------------------------------------------------------
\subsection*{Anexo 2: Matriz de Consistencia Cualitativa}
\addcontentsline{toc}{subsection}{Anexo 2: Matriz de Consistencia Cualitativa}

\begin{table}[H]
\centering
\scriptsize
\begin{tabularx}{\textwidth}{p{2.6cm}p{2.6cm}p{2.6cm}p{2.6cm}X}
\toprule
\textbf{Problema} & \textbf{Objetivos} & \textbf{Categorías} & \textbf{Metodología} & \textbf{Instrumentos} \\
\midrule
\textbf{General:}\newline ¿Cuáles son las vivencias y significados que atribuyen los usuarios al acceso a los servicios de salud y calidad de atención en Ayacucho 2026? & \textbf{General:}\newline Comprender las vivencias y significados sobre el acceso a los servicios de salud y calidad de atención percibida por los usuarios. & \textbf{Categoría 1:}\newline Acceso con Enfoque Intercultural.\newline\newline \textbf{Categoría 2:}\newline Calidad de Atención Percibida. & \textbf{Paradigma:}\newline Interpretativo / Hermenéutico.\newline\newline \textbf{Enfoque:}\newline Cualitativo.\newline\newline \textbf{Método:}\newline Fenomenológico (Epojé). & Guía de Entrevista en Profundidad Bilingüe (8 preguntas Quechua/Español).\newline\newline Diario de campo reflexivo y grabadora de voz. \\
\bottomrule
\end{tabularx}
\caption{Matriz de Consistencia Metodológica Cualitativa.}
\label{tab:consistencia}
\end{table}

\vspace{0.5cm}

% ------------------------------------------------------------------------------
% ANEXO 3: FORMATO DE JUICIO DE EXPERTOS
% ------------------------------------------------------------------------------
\subsection*{Anexo 3: Formato de Validación por Juicio de Expertos}
\addcontentsline{toc}{subsection}{Anexo 3: Juicio de Expertos}

\noindent\textbf{Denominación del Instrumento:} Guía de Entrevista en Profundidad Bilingüe (Quechua Chanka / Castellano).\\
\textbf{Autora:} Gresly Lucero Pariona Palomino \quad | \quad \textbf{Asesor:} Dr. Manglio Aguirre Andrade.

\begin{table}[H]
\centering
\small
\begin{tabularx}{\textwidth}{p{4.0cm}Xcc}
\toprule
\textbf{Criterios de Evaluación} & \textbf{Indicadores Metodológicos} & \textbf{Sí} & \textbf{No} \\
\midrule
1. Claridad & La redacción de las preguntas es comprensible y sin ambigüedades técnicas. & [ X ] & [ \phantom{X} ] \\
2. Objetividad & No induce sesgos ni respuestas preconcebidas en el usuario. & [ X ] & [ \phantom{X} ] \\
3. Pertinencia & Responde con precisión a los objetivos y categorías apriorísticas. & [ X ] & [ \phantom{X} ] \\
4. Adecuación Cultural & La formulación en Quechua respeta la variante Chanka y la cosmovisión andina. & [ X ] & [ \phantom{X} ] \\
5. Consistencia & La estructura fenomenológica garantiza profundidad y saturación teórica. & [ X ] & [ \phantom{X} ] \\
\bottomrule
\end{tabularx}
\caption{Criterios de validación por juicio de expertos.}
\label{tab:expertos}
\end{table}

\vspace{0.5cm}

% ------------------------------------------------------------------------------
% ANEXO 4: CONSENTIMIENTO INFORMADO BILINGÜE
% ------------------------------------------------------------------------------
\subsection*{Anexo 4: Modelo de Consentimiento Informado Bilingüe}
\addcontentsline{toc}{subsection}{Anexo 4: Consentimiento Informado Bilingüe}
\noindent Yo, \underline{\hspace{6cm}}, identificado(a) con DNI N$^\circ$ \underline{\hspace{3cm}}, declaro haber sido informado(a) detalladamente en mi lengua originaria (Quechua Chanka / Castellano) sobre los objetivos del estudio \textit{``Acceso a los servicios de salud con enfoque intercultural y calidad de atención percibida por los usuarios, Ayacucho 2026''}. Acepto voluntariamente participar en la entrevista en profundidad y autorizo la grabación de audio, garantizándose la reserva absoluta, anonimato, confidencialidad estricta y el derecho a retirarme en cualquier momento sin perjuicio de mi atención sanitaria.\\[0.6cm]
\noindent Ayacucho, \underline{\hspace{2cm}} de \underline{\hspace{3cm}} de 2026.\\[0.8cm]
\noindent\begin{tabularx}{\textwidth}{XcX}
\cline{1-1} \cline{3-3}
Firma / Huella Digital del Participante & & Firma de la Investigadora Principal
\end{tabularx}

\vspace{0.5cm}

% ------------------------------------------------------------------------------
% ANEXO 5: CARTA DE ACEPTACIÓN DE ASESORÍA
% ------------------------------------------------------------------------------
\subsection*{Anexo 5: Modelo de Carta de Asesoría de Tesis}
\addcontentsline{toc}{subsection}{Anexo 5: Carta de Asesoría de Tesis}

\noindent Ayacucho, 15 de enero de 2026.\\[0.3cm]
\noindent\textbf{Dr. Alejandro Yarlequé Mujica}\\
Decano de la Facultad de Ciencias de la Salud\\
Universidad Nacional de San Cristóbal de Huamanga\\[0.3cm]
\textbf{Asunto:} Aceptación de Asesoría de Tesis de Grado.\\[0.3cm]
De mi mayor consideración:\\[0.2cm]
Sirva la presente para saludarle muy cordialmente y a la vez comunicarle la \textbf{ACEPTACIÓN FORMAL DE LA ASESORÍA DE TESIS} titulada: \textit{``Acceso a los servicios de salud con enfoque intercultural y calidad de atención percibida por los usuarios, Ayacucho 2026''}, correspondiente a la estudiante de la Escuela Profesional de Enfermería, \textbf{Gresly Lucero Pariona Palomino}.\\[0.2cm]
El plan de tesis cumple plenamente con los lineamientos epistemológicos, metodológicos y éticos de la cátedra de Proyecto de Investigación en Salud (EN 486) y los estándares de la Facultad.\\[0.6cm]
Atentamente,\\[1.0cm]
\noindent\begin{tabularx}{\textwidth}{X}
\centering
\rule{8cm}{0.5pt}\\
\textbf{Dr. MANGLIO AGUIRRE ANDRADE}\\
Docente Titular / Asesor de Tesis\\
Facultad de Ciencias de la Salud -- UNSCH
\end{tabularx}
